The case $p=2$ would give $q^2+r^2=4$, impossible. Thus $p$ is odd and $q,r$ have the same parity. They cannot both be $2$, since $p^2(p-1)\ge18>8$, so all three are odd. Modulo $8$ this gives $p\equiv3\pmod8$.

If none of the primes were $3$, the left side would be $0$ modulo $3$ and the right side nonzero, impossible. If $p=3$, then $q^2+r^2=18$, forcing $q=r=3$.

Suppose instead $p\ne3$ and, after possibly swapping $q,r$, $q=3$. Modulo $p$, the equation gives $r^2\equiv-9$. Since $p\ne3$, this produces a square root of $-1$ modulo $p$ and therefore an element of order $4$. Fermat's theorem implies $4\mid p-1$, contradicting $p\equiv3\pmod8$.

The only triple is $(3,3,3)$, so the average is $\boxed{27}$.
